\section{What remains}
\label{sec:further}

Four things are unfinished. Each is a measurement that can be named, rather
than a direction.

\paragraph{How often the name resolves.}
The condition that decided the live experiment turns on a role's name, and
\nManagedPolicyConfigs deployments attach the policy carrying it. How many of
those roles end up named inside the pattern is unknown, because deployments
commonly build the name from a variable the repository never binds. Counting
that, over the deployments which already attach the policy, would turn the
central result from a mechanism shown in one account into a rate over public
code. It is the highest value item here and it is days of work.

\paragraph{A false alarm rate for detection.}
Section~\ref{sec:defences} measures the first defence and argues the second.
Comparing running code against published code is the only signal that survives
eviction, and how often it would fire on a deployment behaving normally is not
reported here. That needs a population observed while nothing is wrong, which
this corpus cannot supply.

\paragraph{Whether public code resembles what organisations run.}
The prevalence argument rests on an assumption that cannot be tested from
outside, and Section~\ref{sec:limitations} states it as the largest limitation.
A few operators running the analysis over their own accounts would settle more
than another thousand repositories.

\paragraph{The question on a platform that hides its agents.}
Section~\ref{sec:foundry} shows it cannot be asked of source where agents are
not deployment-time resources. Asking it there means reading a running account,
which is a different method against a different population, and it is the
natural next platform rather than a larger corpus.
